\documentclass[11pt]{article}
% ---- style ----
\usepackage[utf8]{inputenc}
\usepackage[T1]{fontenc}
\usepackage{lmodern}
\usepackage{microtype}
\usepackage{amssymb}
\usepackage[a4paper,margin=0.85in]{geometry}
\usepackage{xcolor}
\definecolor{Accent}{HTML}{1565C0}
\usepackage{titlesec}
\titleformat{\section}{\normalsize\bfseries\color{Accent}}{}{0em}{}
\titlespacing*{\section}{0pt}{6pt}{2pt}
\usepackage{enumitem}
\setlist{nosep,leftmargin=1.2em}
\usepackage[colorlinks=true,urlcolor=Accent,linkcolor=Accent]{hyperref}
\input{glyphtounicode}\pdfgentounicode=1
\setlength{\parindent}{0pt}
\setlength{\parskip}{4pt}
\pagestyle{plain}
% ---------------
\begin{document}

\begin{center}
{\large\bfseries\color{Accent} Geometry, Magnetic Fields and Memory of Accreting Neutron Stars:\\
from IXPE to Soft X-ray Polarimetry with GOSoX}\\[3pt]
Research proposal --- Alexander Salganik
\end{center}

\section*{1. Motivation and central question}
Accreting X-ray pulsars (XRPs) are the only systems in which we can watch a
$10^{12}$--$10^{13}$\,G magnetic field channel plasma onto a compact object over five
orders of magnitude in luminosity. Their spectra, pulse profiles and variability encode
where the radiation is produced, how the magnetic field is configured, and how the
accretion disc couples to the magnetosphere. For decades these quantities remained
degenerate. Three tools now break the degeneracy: cyclotron resonant scattering
features (CRSFs) measure the \emph{strength} of the field; IXPE polarimetry measures its
\emph{orientation}, since the polarization angle traces the projected magnetic axis
(e.g.\ Doroshenko et al.\ 2022); and dense monitoring reveals how the flow
\emph{evolves}. In my recent work I discovered that this evolution has a memory: in
1A~1118$-$61 the quasi-periodic oscillation (QPO) frequency at a given luminosity is
$\sim$35\% higher on the rise of the outburst than on the decline. I propose to answer:

\textit{How do the emission geometry, the magnetic-field configuration and the
disc--magnetosphere coupling of accreting neutron stars depend on accretion rate and
accretion history --- and what will soft X-ray polarimetry add?}

\section*{2. Project A --- Geometry and magnetic field from polarimetry and cyclotron lines}
IXPE has observed more than a dozen XRPs. The polarization angle generally follows the
rotating vector model (RVM), giving the first direct measurements of magnetic
obliquities, but the polarization degree is far below theoretical expectations
(Doroshenko et al.\ 2022), and its origin is debated. Because each source has been
analysed with different models and data selections, population questions cannot yet be
asked. I will carry out the first \emph{uniform} spectro-polarimetric analysis of the
IXPE pulsar sample, including data from my own programmes (4U~1954$+$319, GX~1$+$4,
1E~1145.1$-$6141), with pulse-phase-resolved RVM fits performed jointly with NuSTAR
broadband spectra. I will ask: (i) does the polarization degree change between the
low-luminosity atmospheric (Tsygankov et al.\ 2019), sub-critical and supercritical
(Basko \& Sunyaev 1976) regimes? (ii) are obliquities randomly distributed, or are
nearly orthogonal rotators common, as found for X~Persei and GRO~J1008$-$57
(Mushtukov et al.\ 2023)?
\par
\textbf{Geometry meets the cyclotron line.} Once polarimetry fixes the inclination and
obliquity, the angle between the line of sight and the magnetic axis is known at every
pulse phase. I will use it to model the pulse-phase dependence of CRSF energy, width and
depth measured with NuSTAR --- including phase-transient lines such as the one I found
in Swift~J1808.4$-$1754 --- and test whether the lines form in the accretion column or
by reflection from the neutron-star surface (Poutanen et al.\ 2013). This turns two
separate diagnostics, field strength and field orientation, into one consistent map of
the emission region.
\textbf{Deliverables:} a uniform catalogue of XRP geometries; joint
polarization--CRSF constraints on line-forming regions; new IXPE/NuSTAR proposals
targeting regime transitions.

\section*{3. Project B --- The memory of the accretion flow}
QPO hysteresis implies that the inner disc radius lags behind changes in the accretion
rate. I will (i) search for it systematically in Be/X-ray pulsar outbursts with archival
NuSTAR, NICER, Insight-HXMT, SRG/ART-XC and Swift data; (ii) relate its amplitude to the
field strength (from CRSFs), spin period and outburst rise and decay times; and (iii)
build a time-dependent model in which the magnetospheric boundary relaxes on a finite
timescale. The model predicts stronger hysteresis in faster-evolving outbursts, which I
will test with new NICER and NuSTAR target-of-opportunity observations --- an approach
that has already given me $\sim$630\,ks of DDT as PI. The same framework will connect
timing hysteresis with spectral changes such as the emergence of the CRSF-bearing
high-energy hump that I traced in 1A~1118$-$61 at $L\sim10^{36}$\,erg\,s$^{-1}$, and with
transitions into low-luminosity states like the one I found in 2RXP~J130159.6$-$635806.
\textbf{Deliverables:} the first sample study of accretion-history effects in XRPs; a
physical model of disc--magnetosphere ``memory''; predictions for transient monitoring.

\section*{4. Project C --- Accreting pulsars in the soft X-ray polarimetry era}
MKI's GOSoX, selected by NASA for launch in 2030, will measure polarization in the
0.2--0.4\,keV band (Marshall et al.\ 2026). Dedicated science studies have so far focused
on blazars and isolated neutron stars (e.g.\ Kelly et al.\ 2025), while accreting pulsars
remain largely unexplored in this band. Yet this band is where many XRPs show a soft excess with $kT\sim0.1$\,keV,
explained in bright sources by reprocessing of the pulsed hard X-rays on the inner
accretion disc (Hickox et al.\ 2004). Its polarization should differ from that of the
hard emission, so soft X-ray polarimetry can locate the disc--magnetosphere boundary
studied in Project~B. Polarization also changes rapidly with energy: in X~Persei the
degree falls from 30\% at 8\,keV to nearly zero at 2\,keV (Mushtukov et al.\ 2023), so
the soft band cannot be extrapolated from IXPE. Using the MKI CPU/GPU cluster, I will
predict 0.2--0.4\,keV polarization for XRPs, combining the geometries from Project~A
with models of disc reprocessing, scattering and magnetized atmospheres, and assess
feasibility given interstellar absorption, which limits the band to a few sources.
Her~X-1 --- bright, weakly absorbed, with a strongly pulsed soft excess and an
IXPE-measured geometry --- is the prime candidate. At the other end of the spectrum, the
polarization should change sharply across the CRSF, where the two polarization modes
interact very differently with the plasma; I will provide these predictions for future
hard X-ray polarimeters. The fellowship years coincide with GOSoX preparation, so this
work can feed directly into its accreting-pulsar science case and target planning.
\textbf{Deliverables:} soft and hard X-ray polarization predictions for XRPs; an
accreting-pulsar science case and feasibility study for GOSoX (and REDSoX).

\section*{5. Synergy, why MKI, timeline and resources}
The projects reinforce one another: geometries from A feed the timing model of B and the
predictions of C, and C tests the disc--magnetosphere picture of B in a new band. MKI
combines leadership in X-ray polarimetry (REDSoX, GOSoX), the Chandra X-ray Center,
strong neutron-star and X-ray timing groups, and dedicated computing.
% TODO: name the people you have contacted, e.g. "I plan to work closely with Dr. X on ..."
\textbf{Year 1:} uniform IXPE analysis (A); archival hysteresis search (B); first soft
X-ray simulations (C). \textbf{Year 2:} geometry and polarization--CRSF papers; timing
model and ToO campaigns; GOSoX feasibility study. \textbf{Year 3:} population
papers, joint interpretation, new observing and mission proposals. Projects A and B use
data already in hand and are low-risk; C is higher-risk and higher-reward. The research
budget will support travel, campaign coordination and publication costs.

\vspace{2pt}
{\footnotesize\textbf{References.}
Basko \& Sunyaev 1976, MNRAS 175, 395;
Doroshenko et al.\ 2022, Nat.\ Astron.\ 6, 1433;
Hickox et al.\ 2004, ApJ 614, 881;
Kelly et al.\ 2025, ApJ 987, 113;
Marshall et al.\ 2026, arXiv:2608.02891;
Mushtukov et al.\ 2023, MNRAS 524, 2004;
Poutanen et al.\ 2013, ApJ 777, 115;
Salganik et al.\ 2022, MNRAS 514, 2707; 2025, A\&A 698, A71; 2026, A\&A 710, A36;
arXiv:2606.16805; arXiv:2611.20091;
Tsygankov et al.\ 2019, MNRAS 483, L144.}

\end{document}
